Parameterized hardware generation often benefits from templating.
PyTV follows this direction by combining Python as the control language and Verilog as the generated target.
Compared with custom DSL-only approaches such as AHDW \cite{zhao2023automatic},
PyTV intentionally keeps Verilog visible and editable in the source template.

The project provides three complementary interfaces:
\begin{itemize}
  \item A Rust crate (\texttt{pytv}) for integration into toolchains and generators.
  \item A command-line binary (\texttt{pytv}) for direct file conversion.
  \item A Python API binding (\texttt{tverilog}) for in-process generation workflows.
\end{itemize}

At a high level, users write a \texttt{.pytv} file that mixes Verilog and Python directives.
PyTV first converts it into a generated Python script, then optionally executes that script to emit:
\begin{itemize}
  \item a rendered Verilog file (\texttt{.v}), and
  \item a module-instantiation metadata file (\texttt{.inst}) when instantiation blocks are used.
\end{itemize}

PyTV is open-source at \url{https://github.com/autohdw/pytv},
published on crates.io at \url{https://crates.io/crates/pytv},
and distributed under the GPL-3.0-or-later license.
